
\textbf{Principaux domaines abordés :} Probabilités

\medskip

%Au basket-ball, il existe deux sortes de tir :
%
%\begin{itemize}
%\item les tirs à deux points.
%
%Ils sont réalisés près du panier et rapportent deux points s'ils sont réussis.
%\item les tirs à trois points.
%
%Ils sont réalisés loin du panier et rapportent trois points s'ils sont réussis.
%\end{itemize}
%
%\smallskip
%
%Stéphanie s'entraîne au tir. On dispose des données suivantes :
%
%\begin{itemize}
%\item[$\bullet~~$] Un quart de ses tirs sont des tirs à deux points. Parmi eux, 60\,\% sont réussis.
%\item[$\bullet~~$] Trois quarts de ses tirs sont des tirs à trois points. Parmi eux, 35\,\% sont réussis.
%\end{itemize}
%
%\medskip

\begin{enumerate}
\item %Stéphanie réalise un tir. 

%On considère les évènements suivants :
%
%\begin{description}
%\item[ ] $D$ : \og Il s'agit d'un tir à deux points \fg.
%\item[ ] $R$ : \og le tir est réussi \fg.
%\end{description}
	\begin{enumerate}
		\item ~%Représenter la situation à l'aide d'un arbre de probabilités.
\begin{center}
\pstree[treemode=R,nodesepA=0pt,nodesepB=2.5pt,treesep = 1cm,levelsep=2.5cm]{\TR{}}
{\pstree{\TR{$D$~}\taput{0,25}}
{\TR{$R$}\taput{0,6}
\TR{$\overline{R}$}\tbput{0,4}
}
\pstree{\TR{$\overline{D}$~}\tbput{0,75}}
{\TR{$R$}\taput{0,35}
\TR{$\overline{R}$}\tbput{0,65}
}}
\end{center}
		\item %Calculer la probabilité $p\left(\overline{D} \cap R\right)$.
$p\left(\overline{D} \cap R\right) = p\left(\overline{D}\right) \times p_{\overline{D}}\left(R \right) = 0,75 \times 0,35 = \np{0,2625}$.

		\item %Démontrer que la probabilité que Stéphanie réussisse un tir est égale à $0,4125$.
On a de même $p(D \cap R) = p(D) \times p_D(R) = 0,25 \times 0,6 = 0,15$.

D'après la loi des probabilités totales : 

$p(R) = p(D \cap R) + p\left(\overline{D} \cap R\right) =   0,15 + \np{0,2625} = \np{0,4125}$.
		\item %Stéphanie réussit un tir. Calculer la probabilité qu'il s'agisse d'un tir à trois points. Arrondir le résultat au centième.
		Il faut trouver $p_R\left(\overline{D}\right) = \dfrac{p\left(R \cap \overline{D}\right)}{p(R)} = \dfrac{\np{0,2625}}{\np{0,4125}}\approx \np{0,6364}$, soit 0,64 au centième près.
	\end{enumerate}	
\item %Stéphanie réalise à présent une série de $10$ tirs à trois points.

%On note $X$ la variable aléatoire qui compte le nombre de tirs réussis.

%On considère que les tirs sont indépendants. On rappelle que la probabilité que Stéphanie réussisse un tir à trois points est égale à $0,35$.
	\begin{enumerate}
		\item %Justifier que $X$ suit une loi binomiale. Préciser ses paramètres.
Les tirs sont indépendants et à chaque tir la probabilité de le réussir est égale à 0,35 : la variable aléatoire $X$ égale au nombre de réussites suit donc une loi binomiale de paramètres $n = 10$ et $p  = 0,35$.
		\item %Calculer l'espérance de $X$. Interpréter le résultat dans le contexte de l'exercice.
		On a $E = np = 10 \times 0,35 = 3,5$ : en moyenne sur $20$ tirs Stéphanie en réussira $7$.
		\item %Déterminer la probabilité que Stéphanie rate $4$ tirs ou plus. Arrondir le résultat au centième.
		La calculatrice donne $p(X \leqslant 6) \approx 0,97$.
		\item %Déterminer la probabilité que Stéphanie rate au plus $4$ tirs. Arrondir le résultat au centième.
		On a de même $p(X \geqslant 6) \approx 0,095$.
	\end{enumerate}
\item %Soit $n$ un entier naturel non nul.

%Stéphanie souhaite réaliser une série de $n$ tirs à trois points.

%On considère que les tirs sont indépendants. On rappelle que la probabilité qu'elle réussisse un tir à trois points est égale à $0,35$.

%Déterminer la valeur minimale de $n$ pour que la probabilité que Stéphanie réussisse au moins un tir parmi les $n$ tirs soit supérieure ou égale à $0,99$.
La probabilité de rater $n$ tirs à 3 points est égale à $0,65^n$, donc elle d'en réussir un est $1 - 0,65^n$.

Il faut donc résoudre dans $\N$, l'inéquation $1 - 0,65^n \geqslant 0,99 \iff 0,01 \geqslant 0,65^n$, soit par croissance de la fonction logarithme népérien:

$\ln 0,01 \geqslant n\ln 0,62 \iff \dfrac{\ln 0,01}{\ln 0,65} \leqslant n$ (car $\ln 0,65 < 0$).

Or $\dfrac{\ln 0,01}{\ln 0,65} \approx 10,69$ ; le plus petit naturel solution est donc 11 : Stéphanie doit tenter 11 tirs.
\end{enumerate}

\bigskip

